\hypertarget{chapter-1-thinking-like-tex}{%
\chapter{Thinking Like TeX}\label{chapter-1-thinking-like-tex}}

Most introductions to LaTeX begin with commands: how to start a section, how to insert a table, how to number an equation. This book begins somewhere else, because commands are the wrong level of description for a superuser. A LaTeX document is not a sequence of instructions given to a word processor; it is the output of a small, precisely specified computation carried out by TeX itself, with LaTeX supplying the vocabulary in which that computation is expressed. Understanding what TeX actually does with a document, rather than what LaTeX appears to do with it, is the difference between a user who works around the engine's behavior by trial and error and one who predicts it.

\hypertarget{boxes-glue-and-penalties}{%
\section{Boxes, Glue, and Penalties}\label{boxes-glue-and-penalties}}

TeX's model of a typeset page is built from three primitive kinds of material: boxes, glue, and penalties. A box is a rectangular region of fixed width, height, and depth, containing either a piece of typeset text, a rule, or other boxes nested inside it. Every character TeX sets is, at some level, wrapped in a box; every word is a horizontal sequence of character boxes; every line is a horizontal box containing those words; every paragraph is a vertical stack of lines, itself a vertical box.

Boxes alone would produce rigid, unadjustable pages, since a box has a single fixed size. TeX solves this with glue: a flexible space with a natural width and the capacity to stretch or shrink by specified amounts under specified conditions. The space between words is glue, not a fixed-width character; when TeX justifies a line, it is not inserting or deleting spaces but instructing the glue between words to stretch or shrink until the line fills its target width. Vertical space between paragraphs, before and after section headings, and around floats is glue for exactly the same reason: it needs to absorb the differences between one page and the next without being edited by hand.

Penalties are TeX's mechanism for expressing a preference without imposing a rule. A penalty is a number attached to a potential break point in a line or a page; large positive penalties discourage TeX from breaking there, large negative penalties encourage it, and a penalty of \(-10000\) or more forces a break outright. The widow and orphan control built into LaTeX's paragraph and page-breaking behavior, along with commands like \texttt{\textbackslash{}nopagebreak} and \texttt{\textbackslash{}penalty}, are all instructions expressed in this currency.

\textbf{Definition.} A \emph{layout object} in TeX is either a box (fixed extent, opaque content), a glob of glue (an extent with stretch and shrink components), or a penalty (a signed integer attached to a candidate break point). A horizontal or vertical list is a sequence of layout objects; typesetting a paragraph or a page is the process of choosing a subset of the penalties in a list to act as actual breaks, subject to the constraint that the glue in each resulting segment can stretch or shrink to the required target size.

This is worth stating plainly because it explains behavior that otherwise looks like a bug. A paragraph that refuses to justify without visibly stretched spacing has run out of glue to stretch; a section heading that jumps to the next page did so because a penalty forbade breaking immediately after it; a table that will not sit where it was written moved because floats are boxes subject to their own placement rules, not because LaTeX is arbitrarily uncooperative. Once boxes, glue, and penalties are visible as the actual material TeX is working with, most of what look like idiosyncrasies of the engine turn out to be the direct, legible consequence of a small set of rules applied consistently.

\hypertarget{tokens-category-codes-and-the-lexerexpander-split}{%
\section{Tokens, Category Codes, and the Lexer/Expander Split}\label{tokens-category-codes-and-the-lexerexpander-split}}

Before any of that layout machinery runs, TeX has to turn a stream of characters in a source file into something it can act on. This happens in two stages that are worth keeping separate in one's head, because most confusing macro behavior comes from conflating them.

The first stage is tokenization. TeX reads raw characters and, using a table of \emph{category codes}, converts each one into a token: a control sequence (\texttt{\textbackslash{}section}, \texttt{\textbackslash{}newcommand}), a character with a category such as letter, space, math shift, or alignment tab, or one of a small number of special tokens. Category codes are not fixed properties of characters; they are assignments that can be changed. The fact that \texttt{\#} means ``parameter marker'' inside a macro definition but a literal hash character when its category code is altered, or that \texttt{\%} can be made to stop being a comment character, follows from the same mechanism: TeX does not know what a character ``means'' until it has consulted the category code table current at the moment that character is read.

The second stage is expansion. An expandable token, such as a macro name, is repeatedly replaced by its definition until only unexpandable tokens remain: primitive commands, characters ready to be typeset, or explicit instructions like box and glue specifications. This is the stage most macro-writing takes place in, and it is where the rewriting-system character of TeX becomes visible. Defining a macro with \texttt{\textbackslash{}newcommand} or \texttt{\textbackslash{}def} does not create a procedure that TeX will later call; it registers a rewrite rule that TeX applies wherever the macro's name appears, exactly as many times as expansion requires and no more.

Keeping tokenization and expansion apart clarifies a class of problems that otherwise look mysterious: a macro that behaves differently depending on what character follows it, a command that fails only inside a \texttt{verbatim} environment, or a package that changes the meaning of an ordinary character such as \texttt{"} or \texttt{\textasciitilde{}}. In each case, the category code table active at the point a character is read has changed, and the tokens produced from the same visual character sequence are simply different tokens than they would otherwise have been.

\hypertarget{expansion-order-what-tex-sees-versus-what-tex-does}{%
\section{Expansion Order: What TeX Sees versus What TeX Does}\label{expansion-order-what-tex-sees-versus-what-tex-does}}

A recurring source of difficulty for LaTeX users moving from ordinary commands to writing their own macros is that expansion does not proceed left to right through a document in the way execution would in a conventional programming language. TeX expands tokens lazily, one at a time, only as they are needed by whatever process is currently consuming input, whether that is the paragraph builder, a conditional, or another macro's argument-gathering.

This produces two behaviors that are easy to misjudge without the model in hand. First, an argument to a macro is not evaluated before the macro is called; it is handed over as an unexpanded token sequence, and whether or how it gets expanded depends entirely on what the macro's definition does with it. A macro that never uses its argument, or uses it inside a context that suppresses expansion such as a section title written to an auxiliary file, may leave nested macros in that argument completely unexpanded, which is why commands like \texttt{\textbackslash{}protect} and the \texttt{\textbackslash{}text}-style variants for headings and captions exist at all: they manage exactly this asymmetry between where a token sequence appears and where it is ultimately consumed.

Second, controlling when expansion happens, rather than only what a macro expands to, is itself a design tool. The primitive \texttt{\textbackslash{}expandafter} lets a macro writer reach past one token to expand the one behind it first, which is the standard technique for building a new control sequence name out of pieces, or for testing the first character of an argument before deciding how to process the rest. Packages built for robust macro programming, such as \texttt{expl3}, replace this style of manual token-shuffling with a more disciplined function-based interface, but the underlying model they sit on top of is the same lazy, rule-based expansion described here.

\textbf{Example.} Consider a macro intended to typeset a term and, the first time it is used, also add that term to a glossary:

\begin{verbatim}
\newcommand{\term}[1]{\textit{#1}\glossaryadd{#1}}
\end{verbatim}

Here \texttt{\#1} is not a value computed once and reused; it is a token sequence substituted verbatim into the macro's replacement text at every point it appears. If the argument itself contains a macro, such as a cross-reference or another custom command, that inner macro will be expanded independently, in place, whenever and wherever the surrounding context causes expansion to reach it, not once at the moment \texttt{\textbackslash{}term} is written.

None of this is presented for its own sake. The chapters that follow, on writing key-value interfaces, designing document classes, and building large multi-file projects, all depend on treating macros as rewrite rules operating on token streams rather than as function calls in a conventional sense. A superuser's fluency with LaTeX comes less from memorizing a larger command vocabulary than from being able to predict, from the model in this chapter, what a given piece of source will actually become once TeX has finished reading it.
